\section{Evidencia de la auditor\'ia pre-campo: los artefactos que quedaron}
\label{sec:evidencia-auditoria}

La auditor\'ia previa a la campa\~na dej\'o algo poco habitual y muy valioso:
\textbf{archivos de evidencia con el antes y el despu\'es de cada correcci\'on},
en formato legible por m\'aquina, en
\texttt{src/firmware/esp32/Nodo comunicaci\'on/master/artifacts/}. No son notas
sobre lo que pas\'o: son el registro que la propia herramienta de prueba escribi\'o
mientras pasaba.

\subsection{El fallo F1 y su correcci\'on, en n\'umeros}

El caso mejor documentado es el defecto \textbf{F1}: el \emph{watchdog} de
\texttt{SYNC\_STALL} destru\'ia las capturas largas a tarjeta SD. Existen los dos
archivos, el de la corrida fallida y el de la corrida que pas\'o, con la misma
configuraci\'on nominal:

\begin{center}\footnotesize
\begin{tabularx}{\textwidth}{@{}l r r X@{}}
\toprule
& \textbf{Antes (falla)} & \textbf{Despu\'es (PASS)} & \\
\midrule
Muestras escritas & \textbf{0} & \textbf{\num{1562400}} & de \num{1562400} esperadas \\
Lotes & 0 & \num{52080} & de 30 muestras cada uno \\
Bytes en disco & 0 & \num{4687200} & $=\num{1562400}\times3$ (entero de 24 bits) \\
Duraci\'on pedida & \SI{600}{\second} & \SI{600}{\second} & a \SI{2604}{\hertz}, sin decimar \\
Tiempo transcurrido & \SI{3420.6}{\second} & \SI{2626.5}{\second} & \\
\textbf{Desconexiones} & \textbf{15} & \textbf{0} & \\
\textbf{Conexiones} & \textbf{16} & \textbf{1} & \\
Limpieza final & fallida & correcta & \\
Estado & \texttt{error} & \texttt{ok} & \\
\bottomrule
\end{tabularx}
\end{center}

El mensaje de error de la corrida fallida se conserva literal y describe el modo
de falla con precisi\'on: \texttt{PhaseTimeout: VER/captura timeout:
muestras=0/1562400}, seguido de \texttt{final\_cleanup/STOP no confirmado: falta
ACK 0xA4}. Las \textbf{15 desconexiones} y \textbf{16 reconexiones} son la firma
del \emph{watchdog} re-disparando: cada reintento cortaba el enlace y reiniciaba
la sesi\'on, de modo que la captura nunca llegaba a completarse.

\seleccionado{\textbf{La aritm\'etica cierra exactamente} en la corrida que pasa:
$\num{52080}$ lotes $\times\,30$ muestras $=\num{1562400}$ muestras, que a
\SI{2604}{\hertz} son exactamente \SI{600}{\second}; y
$\num{1562400}\times3$ bytes $=\num{4687200}$ bytes, coincidentes con el tama\~no
real del archivo. Ni una muestra de m\'as (\texttt{extra\_samples}${}=0$) ni de
menos. \textbf{Cero desconexiones y una sola conexi\'on} en 44 minutos de
operaci\'on continua.}

\dato{\textbf{Un dato de rendimiento que conviene tener presente}: capturar
\SI{600}{\second} llev\'o \SI{2626}{\second} de reloj, es decir que la
\textbf{transferencia de los datos tarda unas cuatro veces m\'as que la captura}.
La captura en s\'i es tiempo real; el volcado por WebSocket con esquema
\emph{stop-and-wait} es el cuello de botella. Para una jornada de campo con
capturas largas esto no es un detalle: define cu\'antos puntos se alcanzan a
medir por d\'ia.}

\subsection{Otros artefactos conservados}

\begin{center}\footnotesize
\begin{tabularx}{\textwidth}{@{}l r X@{}}
\toprule
Artefacto & Muestras & Qu\'e demuestra \\
\midrule
\texttt{geo\_sd\_smoke\_5} & --- & Prueba de humo m\'inima de la SD (3 variantes: base, \texttt{\_v2} y \texttt{\_pass}) \\
\texttt{geo\_sd\_medium\_600} & \num{18000} & 600 lotes, \SI{54000}{} bytes, 0 desconexiones, \SI{34.6}{\second} \\
\texttt{e5c\_geo10min\_2604\_pass} & \num{1562400} & La corrida de \SI{10}{\minute} completa que pasa \\
\texttt{geo\_10min\_2604} & 0 & La corrida fallida, conservada como evidencia del defecto \\
\texttt{f9\_start\_flow\_pass} & --- & Los escenarios del flujo de START tras corregir F9 \\
\bottomrule
\end{tabularx}
\end{center}

\seleccionado{\textbf{Conservar el archivo de la corrida fallida es la decisi\'on
correcta} y merece se\~nalarse. Lo natural habr\'ia sido borrarlo al conseguir el
PASS. Mantenerlo permite hoy, un mes despu\'es, reconstruir exactamente qu\'e
fallaba y cu\'anto mejor\'o ---y es lo que hace posible la tabla de arriba.}

\subsection{Los barridos crudos de osciloscopio}

En \texttt{data/raw/Osciloscopio/} se conservan los archivos originales del
osciloscopio, organizados \textbf{una carpeta por d\'ecada de barrido}:

\begin{center}\footnotesize
\begin{tabular}{@{}ll@{}}
\toprule
\multicolumn{2}{@{}l}{15 carpetas, 4 archivos CSV cada una (un canal por archivo) $=$ \textbf{60 CSV crudos}} \\
\midrule
\texttt{100mHz\_1Hz} & \texttt{100mHzto1Hz\_200s\_25s\_div} \\
\texttt{500mHz\_5Hz} & \texttt{500mHzto5Hz\_40s\_5s\_div} \\
\texttt{2.5Hz\_25Hz} & \texttt{2.5Hzto25Hz\_8s\_1s\_div} \\
\texttt{10Hz\_100Hz} & \texttt{10Hzto100Hz\_2s\_250ms\_div} \\
\texttt{30Hz\_300Hz} & \texttt{50Hzto500Hz\_400ms\_50ms\_div} \\
\texttt{100Hz\_1kHz} & \texttt{250Hzto2.5kHz\_80ms\_10ms\_div} \\
--- & \texttt{1kHzto10kHz\_20ms\_2.5ms\_div} \\
--- & \texttt{5kHzto50kHz\_4ms\_500us\_div} \\
--- & \texttt{20kHzto200kHz\_1ms\_100us\_div} \\
\bottomrule
\end{tabular}
\end{center}

La nomenclatura registra la banda \emph{y} la base de tiempo del osciloscopio, de
modo que cada archivo es reproducible sin consultar una libreta. Las dos
convenciones de nombre (\texttt{100mHz\_1Hz} frente a
\texttt{100mHzto1Hz\_200s\_25s\_div}) corresponden a dos tandas distintas; la
segunda, m\'as expl\'icita, es la que se mantuvo.

\subsection{La m\'aquina de estados de la calibraci\'on}

\begin{figure}[ht]
\centering
\includegraphics[width=0.98\textwidth]{mosaico_fsm_calibracion.png}
\caption{Las dos versiones del diagrama de la m\'aquina de estados de
auto-calibraci\'on, generadas por \texttt{docs/calibration\_fsm/draw\_calibration\_fsm.py}:
la simplificada, que es la que se us\'o en el art\'iculo, y la detallada, con
todas las transiciones y guardas.}
\label{fig:mosaico-fsm}
\end{figure}

\seleccionado{El diagrama est\'a \textbf{generado por c\'odigo}
(\texttt{draw\_calibration\_fsm.py}) y no dibujado a mano, lo que significa que
puede regenerarse si la m\'aquina cambia. Acompa\~na un
\texttt{calibration\_fsm\_explanation.md} con la descripci\'on de cada estado.}